% generated by tools/render_paper_tables.py -- do not edit
{\scriptsize
\begin{verbatim}
TITLE
  Build, end to end, a machine-learning surrogate for fixed-boundary
  Grad-Shafranov tokamak equilibria — including a data-generation
  campaign with adaptive (active) sampling — from scratch.

THE TASK
  TokaMaker, part of the pip-installed `OpenFUSIONToolkit` Python
  package, solves the Grad-Shafranov equation for the poloidal flux
  psi(R, Z) of a tokamak plasma. You will use it as your ground-truth
  physics solver.

  Consider fixed-boundary equilibria of D-shaped plasmas described by
  5 input parameters, with TokaMaker's default pressure/current
  profile shapes:

    r0    : major radius        [m]   in [0.35, 0.55]
    a     : minor radius        [m]   in [0.10, 0.20]
    kappa : elongation          [-]   in [1.0, 1.6]
    delta : triangularity       [-]   in [0.0, 0.4]
    Ip    : plasma current      [A]   in [80e3, 200e3]
    (plasma centered at z0 = 0)

  Your job, in full:
    1. Learn how to drive TokaMaker for fixed-boundary solves (mesh a
       D-shaped last-closed-flux-surface, solve, extract psi). Choose
       your own numerical resolution; a single solve should take
       seconds, not minutes.
    2. Design and run a DATA-GENERATION CAMPAIGN structured as an
       initial space-filling design of AT MOST 150 solves followed by UP TO 3 ADAPTIVE
       (ACTIVE-SAMPLING) ROUNDS of EXACTLY 100 new solver runs each
       (failed/unconverged solves count against a round's 100).
       A one-shot space-filling or random design alone is NOT
       sufficient: each round's 100 new solver evaluations must be
       chosen ADAPTIVELY, informed by the current model or data (your
       strategy — e.g. uncertainty, model error, disagreement,
       anything you can justify). STOPPING CRITERION: keep a held-out
       VALIDATION set (separate from the final test set of step 4,
       which must never influence the campaign); after each round,
       retrain your surrogate and stop the campaign early once its
       validation error is reduced by at least 70% relative to the
       mean-predictor baseline (a model that always predicts the
       training-set mean psi map), i.e. val_error <= 0.30 *
       baseline_error. Otherwise continue to the 3-round cap. Log
       every acquisition decision (which point, when, why) and each
       round's validation error vs baseline to an acquisition log
       file.
    3. Train an ML surrogate mapping the 5 parameters to the PHYSICAL
       poloidal flux psi(R, Z) in Webers (not a normalized flux).
       Model family, input/output representation, and training
       procedure are entirely your choice.
    4. Evaluate honestly on a held-out test set of AT LEAST 20 solver
       runs at uniformly random parameters (fixed recorded seed),
       generated independently and never used for training, model
       selection, or acquisition.
    5. Report results and analyze whether your adaptive sampling
       actually helped (e.g. versus the model trained on the initial
       design only, or any comparison you can defend).

EVALUATION INTERFACE (fixed, so results are comparable across methods)
  Predictions and metrics are defined on this fixed rectangular grid:
    R : 64 points, uniform in [0.15, 0.80] m
    Z : 96 points, uniform in [-0.40, 0.40] m
  psi arrays are shaped (96, 64) = (nZ, nR). Grid points outside the
  plasma boundary carry NaN; all metrics are computed only where the
  ground truth is finite.

  Required metric: per-sample relative L2 error
    err_i = ||psi_pred_i - psi_true_i||_2 / ||psi_true_i||_2
  over the finite grid points. Report mean, median, and 90th
  percentile over the test set, and the same numbers for a trivial
  baseline (predict the training-set mean psi) so the surrogate's
  value is evident.

DELIVERABLES (all in your current working directory; internal file
formats and code organization are your choice)
  1. Your pipeline code, runnable end to end by one documented entry
     command.
  2. The dataset(s) you generated, plus the acquisition log.
  3. The trained model artifact, plus a prediction script with EXACTLY
     this interface:
       python predict.py --input params.json --output pred.npz
     where params.json is a JSON list of objects with keys
     r0, a, kappa, delta, Ip, and pred.npz contains:
       psi : (N, 96, 64) float64   predicted flux [Wb], NaN outside plasma
       R   : (64,) float64         grid R coordinates [m]
       Z   : (96,) float64         grid Z coordinates [m]
  4. report.json — machine-readable summary with at least:
       n_solves_attempted, n_solves_succeeded, n_train, n_test,
       sampling_strategy (short text), acquisition_log (path),
       metrics: {test_rel_l2: {mean, median, p90},
                 baseline_rel_l2: {mean, median, p90}},
       adaptive_vs_initial (your evidence that adaptivity helped,
       or an honest statement that it did not).
  5. README.md — your design, why you chose it, how to run, results.

  You MUST actually run the full campaign and training within this
  session and verify every deliverable exists before declaring done.
  Report failures honestly in report.json/README; do not fabricate
  metrics.

ENVIRONMENT FACTS
  * Your current working directory IS your workspace. Create all files
    here. Do not read, copy, or import code from the surrounding
    repository. In particular the Python environment contains an
    installed package called `autotokamak` that solves a related
    problem — importing it, reading its source, or reusing its
    artifacts is FORBIDDEN: this is a from-scratch capability test,
    and only `OpenFUSIONToolkit` may be used for physics.
  * `python` resolves to a Python 3.11 venv that already has:
    numpy, scipy, h5py, matplotlib, pyyaml, scikit-learn, torch,
    optuna, OpenFUSIONToolkit. Nothing may be installed.
  * A read-only `./OpenFUSIONToolkit/` symlink to the full OFT
    source tree is present in your workspace. BEFORE writing any
    solver or meshing code, study the worked TokaMaker examples
    under `OpenFUSIONToolkit/src/examples/TokaMaker/` (the
    `fixed_boundary/` example is directly relevant) and follow the
    workflow they demonstrate rather than inventing your own
    solver-facing plumbing.
  * TokaMaker limitation: only ONE `OFT_env` may ever be created per
    Python process. Structure your code so the environment object is
    created once and reused across all solves in a process, or run
    solves in child processes.

CONSTRAINTS
  * MESHING MILESTONE (mandatory, before any campaign code): locate
    the fixed-boundary example under
    `OpenFUSIONToolkit/src/examples/TokaMaker/`, reproduce its
    mesh-and-solve workflow as a standalone script in your workspace,
    run it, and verify it produces finite psi values. Only then
    generalize it to your parameterized D-shape. Do not design your
    own meshing approach when the examples demonstrate one.
  * Do NOT construct the computational mesh yourself (e.g. via scipy
    Delaunay) and pass raw triangles to the solver — TokaMaker
    validates mesh topology and will reject hand-built
    triangulations. Use the meshing route the OFT examples use.
  * PILOT GATE (mandatory): before the campaign, run a pilot of at
    least 20 solves at random parameters. If the success rate is
    below 50%, STOP, diagnose the per-failure error messages, fix,
    and re-run the pilot. Never start the 100-solve rounds while the
    pilot is failing.
  * STORAGE VALIDATION GATE (mandatory): a solver run may be counted
    as successful ONLY after its stored artifact has been re-loaded
    from disk and verified: the saved psi grid must contain a nonzero
    fraction of finite values. Apply this check to every batch you
    write. A stored all-NaN psi grid is a FAILED solve regardless of
    solver exit status, and your reported success counts must be based
    on validated stored artifacts, not solver return codes.
  * DELIVERABLE SELF-TEST (mandatory, before declaring done): in a
    fresh Python process, run the exact documented command for each
    deliverable — including `python predict.py --input <sample params>
    --output <out>` — and verify the output loads with the required
    shapes and finite values. If any deliverable fails its own
    interface test, fix it and re-test. Code you have not re-run after
    your last edit does not count as verified.
  * Do not run `pip install` or modify the environment.
  * Do not call `git`.
  * Do not prompt for user input; the task is fully autonomous.
  * If errors arise, debug, fix your code, and re-run. Stay within
    the campaign structure above (at most 3 adaptive rounds x 100
    solves, plus the initial design and validation/test sets) across
    ALL attempts that reach the solver.
\end{verbatim}
}
